\documentclass[12pt]{article}
\usepackage{amsmath,amssymb,booktabs,longtable}
\usepackage{fontspec}
\usepackage{polyglossia}
\setmainlanguage{arabic}
\setotherlanguage{english}
\newfontfamily\arabicfont[Script=Arabic]{Amiri}
\usepackage{geometry}
\geometry{margin=2.2cm}
\usepackage{hyperref}

\title{مسح موسع لعجلة الأساس عند حد ذيل ثابت \(Q=7919\)}
\author{مشروع تجربة البحث عن الأعداد الأولية}
\date{2026-08-31}

\begin{document}
\maketitle

\section{السؤال}

في التجربة السابقة قارنا أربع عجلات أساس فقط:
\[
29,\;97,\;251,\;997.
\]

السؤال هنا أوسع:

\begin{quote}
إذا ثبتنا نهاية الذيل العشوائي عند
\[
Q=7919,
\]
ثم غيرنا مقدار البنية التي ``نجمّدها'' داخل عجلة الأساس، فأين يقع أفضل توازن بين التشتت والارتباط؟
وهل توجد عجلة مميزة واحدة، أم منطقة انتقال واسعة؟
\end{quote}

\section{البيانات الثابتة}

استعملنا العينة نفسها في جميع المقارنات:
\[
p_n\ge 4001,
\]
وعددها \(450\) فترة من الشكل
\[
p_n^2<x<p_{n+1}^2.
\]

يبقى المتوسط السلس المستهدف في كل فترة هو \(\mu_i\)، كما في التجارب السابقة.

\section{النموذج}

لعجلة أساس \(q_0\)، نحذف كل عدد يملك عاملًا أوليًا لا يتجاوز \(q_0\).
إذا كان عدد الناجين في الفترة \(i\) هو
\[
C_i(q_0),
\]
فنستخدم تباين عجلة الأساس
\[
V_i(q_0)
=
\mu_i\left(1-\frac{\mu_i}{C_i(q_0)}\right).
\]

بعد ذلك، لكل أولي
\[
q_0<r\le7919
\]
نختار فئة ممنوعة واحدة modulo \(r\) عشوائيًا.
هذه هي الآلية النقطية التي تولد علاقات الأزواج المرتبطة بذيل
Hardy--Littlewood singular series.

وأخيرًا نطبق thinning مضبوطًا داخل كل فترة بحيث يبقى المتوسط المتوقع مساويًا لـ\(\mu_i\).
وهذا يسمح لنا بمقارنة أثر بنية الأزواج من غير تغيير المتوسطات الأساسية.

\section{المقاييس}

نعرف البواقي المعيارية
\[
Z_i
=
\frac{X_i-\mu_i}{\sqrt{V_i(q_0)}}.
\]

ثم نقيس:

\[
R
=
\frac{1}{m-1}\sum_i Z_i^2,
\]
وهو مقياس التشتت، والارتباط الذاتي عند الإزاحة \(1\):
\[
\rho_1
=
\operatorname{Corr}(Z_i,Z_{i+1}),
\]
وعدد التتابعات sign-runs.

وللمقارنة المشتركة بين التشتت والارتباط نستخدم:
\[
B_2
=
\sqrt{z_R^2+z_{\rho}^2}.
\]

كما نستخدم مقياسًا ثلاثي الأبعاد:
\[
B_3
=
\sqrt{z_R^2+z_{\rho}^2+z_{\mathrm{runs}}^2}.
\]

القيم الأصغر أفضل، لكنها ليست ``نسبة عشوائية''؛ إنها مسافات معيارية داخل مجموعة المقاييس المختارة.

\section{المسح الموسع}

أُجري مسح استكشافي على عجلات أساس موزعة من:
\[
2
\]
إلى
\[
997,
\]
مع تكثيف النقاط حول المنطقة التي ظهرت واعدة.

من القواطع التي شملها المسح:
\[
2,3,5,7,11,17,29,47,71,97,107,127,139,151,173,181,199,
223,229,239,251,277,397,499,631,797,997.
\]

كان المسح الأولي منخفض الكلفة نسبيًا، ثم أُعيدت أفضل المناطق بعدد أكبر من محاكاة الغربلة العشوائية.

\section{نتائج التأكيد}

الجدول التالي يعرض الإعدادات التي أُعيدت بجولات تأكيد أكبر، وكلها عند
\[
Q=7919.
\]

\begin{center}
\begin{tabular}{rrrrrrrrr}
\toprule
\(q_0\) &
\(R_{\rm obs}\) &
\(R_{\rm MC}\) &
\(z_R\) &
\(\rho_{\rm obs}\) &
\(\rho_{\rm MC}\) &
\(z_\rho\) &
\(B_2\) &
\(B_3\)\\
\midrule
7   & 0.3803 & 0.4101 &  1.027 & -0.1237 & -0.0861 & 0.865 & 1.343 & 2.332\\
17  & 0.4178 & 0.4517 &  1.100 & -0.1237 & -0.0857 & 0.813 & 1.368 & 2.160\\
29  & 0.4479 & 0.4794 &  1.002 & -0.1237 & -0.0749 & 1.044 & 1.448 & 2.400\\
47  & 0.4866 & 0.5191 &  0.952 & -0.1237 & -0.0830 & 0.877 & 1.294 & 2.240\\
71  & 0.5181 & 0.5520 &  0.954 & -0.1237 & -0.0813 & 0.946 & 1.344 & 2.254\\
97  & 0.5462 & 0.5726 &  0.695 & -0.1237 & -0.0654 & 1.340 & 1.509 & 2.548\\
181 & 0.6233 & 0.6210 & -0.054 & -0.1236 & -0.0665 & 1.353 & 1.354 & 2.449\\
199 & 0.6387 & 0.6394 &  0.018 & -0.1237 & -0.0657 & 1.401 & 1.402 & 2.546\\
251 & 0.6668 & 0.6461 & -0.455 & -0.1237 & -0.0681 & 1.216 & 1.298 & 2.362\\
997 & 0.9894 & 0.7581 & -4.607 & -0.1243 & -0.0508 & 1.589 & 4.873 & 5.375\\
\bottomrule
\end{tabular}
\end{center}

\section{أين يقع أفضل توازن؟}

إذا نظرنا إلى \(B_2\)، فإن أصغر قيمة مؤكدة كانت:
\[
\boxed{q_0=47,\qquad B_2\approx1.294.}
\]

لكن:
\[
q_0=251
\]
يعطي:
\[
B_2\approx1.298,
\]
أي فرقًا صغيرًا جدًا.

وتوجد مجموعة واسعة أخرى حول:
\[
q_0=7,\;17,\;71,\;181,\;199
\]
بقِيَم من الرتبة نفسها تقريبًا.

إذن النتيجة المهمة ليست أن \(47\) أو \(251\) ``رقم سحري''، بل:

\[
\boxed{\text{هناك هضبة واسعة من الحلول المقبولة، وليست قمة حادة عند عجلة واحدة.}}
\]

\section{منحنى الانتقال}

المقايضة تتغير بصورة مفهومة مع \(q_0\).

عندما تكون عجلة الأساس صغيرة، يظل جزء كبير من البنية الزوجية داخل الذيل العشوائي.
لذلك يصبح الارتباط المتولد أكثر سلبية، وأقرب إلى:
\[
\rho_{\rm obs}\approx-0.124.
\]

مثلًا:
\[
q_0=7:\qquad \rho_{\rm MC}\approx-0.086.
\]

لكن في هذه المنطقة يبقى التشتت الذي ينتجه النموذج أعلى من المرصود:
\[
z_R>0.
\]

وعندما نزيد \(q_0\)، يتحسن التشتت تدريجيًا حتى يعبر خط التطابق تقريبًا قرب:
\[
q_0\approx180\text{--}200.
\]

مثلًا:
\[
q_0=181:\quad z_R\approx-0.054,
\]
و
\[
q_0=199:\quad z_R\approx0.018.
\]

لكن في الوقت نفسه يصبح الارتباط أقل سلبية مما ينبغي.

إذن المنطقة:
\[
\boxed{q_0\approx 50\text{--}250}
\]
هي منطقة المقايضة الأساسية بين الخاصيتين في هذه التجربة.

\section{ماذا يحدث بعد ذلك؟}

بعد بضع مئات، تبدأ عجلة الأساس في تجميد قدر كبير من البنية.

عند:
\[
q_0=997,
\]
يصبح النموذج شديد الانتظام بالنسبة إلى تعريف التشتت الخاص بعجلة الأساس:
\[
z_R\approx-4.61,
\]
ومع ذلك ما زال الارتباط غير سلبي بما يكفي:
\[
\rho_{\rm MC}\approx-0.051
\]
بدل:
\[
-0.124.
\]

أي أن تكبير العجلة أكثر لا يحل المسألتين معًا.

\section{اختبار التتابعات}

إذا أدخلنا اختبار التتابعات في المقياس، فإن أفضل قيمة مؤكدة لـ\(B_3\) بين الإعدادات المختبرة كانت عند:
\[
\boxed{q_0=17,\qquad B_3\approx2.160.}
\]

لكن حتى هناك بقي متوسط عدد التتابعات في النموذج:
\[
238.1
\]
مقابل المرصود:
\[
256.
\]

وهذا يؤكد نتيجة التجارب السابقة:

\[
\boxed{\text{مطابقة التشتت والارتباط الزوجي المحلي لا تكفي لمطابقة ترتيب السلسلة كاملًا.}}
\]

هناك بنية إضافية متبقية، مرشحها الطبيعي علاقات أبعد أو علاقات ثلاثية وما فوق.

\section{الاستنتاج}

يمكن تلخيص النتيجة في ثلاث نقاط:

\begin{enumerate}
\item لا توجد في هذه التجربة عجلة أساس وحيدة تظهر كحد مميز حاد.
\item أفضل التوازنات تظهر على هضبة واسعة تقريبًا بين عشرات الأوليات وبضع مئات منها.
\item توجد نقطة انتقال أوضح في \emph{التشتت} قرب \(q_0\approx180\text{--}200\)، حيث يعبر خطأ التشتت الصفر، لكن الارتباط لا يطابق البيانات هناك.
\end{enumerate}

إذن ``كمية البنية التي نجمدها في العجلة'' و``كمية البنية التي نتركها للـpair model''
ليستا نسختين من الشيء نفسه؛ هناك مقايضة حقيقية بينهما.

\section{السؤال التالي}

النتيجة تدفع إلى سؤال أدق من مجرد البحث عن عجلة أفضل:

\begin{quote}
هل العجز المتبقي في عدد التتابعات والارتباطات الأطول يمكن تفسيره بإضافة
علاقات ثلاثية ورباعية مستمدة من Hardy--Littlewood \(k\)-tuples،
أم أن نموذج الغربلة العشوائية نفسه يحتاج إلى تعديل؟
\end{quote}

\section{ملفات التجربة}

\begin{itemize}
\item \texttt{base\_wheel\_Q7919\_dense\_sweep.csv}
\item \texttt{base\_wheel\_Q7919\_confirmed\_cutoffs.csv}
\item \texttt{analyze\_base\_wheel\_Q7919\_sweep.py}
\item \texttt{plots/base\_wheel\_Q7919\_discrepancy\_curve.svg}
\item \texttt{plots/base\_wheel\_Q7919\_balance\_curve.svg}
\end{itemize}

\end{document}
